% Algebra 2 Lesson 2-7 fast-track packet: Linear-Quadratic Systems (enVision A2 2024 TE, Example 3 and Example 4 only).
% Fast-lesson draft 2026-10-30 (docs/a2-fast-lesson-spec.md); PLC has not met on this lesson.
% Quiz map (TE 2-7 Lesson Quiz): item 1 count solutions = Ex 1; item 2 substitution = Ex 2; item 3 ball up a hill = Ex 3;
% item 4 shade a system of inequalities = Ex 4; item 5 calculator intersect = Ex 5.
% Kept: Ex 3 (the lesson's only DOK 3, practice 29; quiz 3, and its set-equal step is all of quiz 2) and
% Ex 4 (the only route to quiz 4). DOK 3: practice 29 is the Practice block (page 2); the Exit question is Additional Practice 5.
% Quiz numbers below are the Savvas item numbers (the short form keeps them); items 1 and 5 are full Savvas quiz only.
% Shared preamble for the "fast" (hacked-down) Algebra 2 lesson packets, Oct 2026.
% Each packet: % Shared preamble for the "fast" (hacked-down) Algebra 2 lesson packets, Oct 2026.
% Each packet: % Shared preamble for the "fast" (hacked-down) Algebra 2 lesson packets, Oct 2026.
% Each packet: \input{fast-packet-template} then \begin{document} ... \end{document}.
\documentclass[12pt]{article}
\usepackage[margin=0.7in]{geometry}
\usepackage{amsmath,amssymb}
\usepackage{booktabs,tabularx,array}
\usepackage{enumitem}
\usepackage[T1]{fontenc}
\usepackage[utf8]{inputenc}
\usepackage{lmodern}
\usepackage{tikz}
\usepackage{pgfplots}
\pgfplotsset{compat=1.18}
\usepackage{xcolor}
\setlist{nosep,leftmargin=1.4em}
\setlength{\parindent}{0pt}
\setlength{\parskip}{6pt}
\renewcommand{\arraystretch}{1.35}
\newcolumntype{Y}{>{\raggedright\arraybackslash}X}
\pagestyle{empty}

% blanks and workspace
\newcommand{\blank}[1][2.6cm]{\rule[-2pt]{#1}{0.5pt}}
\newcommand{\lines}[1][1]{\par\vspace{2pt}\foreach \i in {1,...,#1}{\rule{0pt}{22pt}\rule{\linewidth}{0.4pt}\par}}
\newcommand{\wspart}[2]{\vspace{4pt}{\large\bfseries #1: #2}\par\vspace{-6pt}\rule{\textwidth}{0.4pt}\par}

% header: \packethead{Lesson 1-3}{Piecewise-Defined Functions}
\newcommand{\packethead}[2]{%
{\LARGE\bfseries #2}\hfill Name \blank[5cm]\\[2pt]
Algebra 2 \quad #1 (fast track) \hfill Date \blank[2.5cm] \quad Period \blank[1.2cm]\par\vspace{2pt}}

% boxed goals: \goals{learning goal}{language goal}
\newcommand{\goals}[2]{%
\fbox{\parbox{\dimexpr\textwidth-2\fboxsep-2\fboxrule}{%
\textbf{Learning goal:} #1\\[2pt]
\textbf{Language goal:} #2}}\par}

% quiz target box: \quiztarget{text}
\newcommand{\quiztarget}[1]{%
\vspace{4pt}\colorbox{gray!12}{\parbox{\dimexpr\textwidth-2\fboxsep}{\textbf{On the quiz:} #1}}\par}

% IXL table row helper: code & skill & which quiz question it feeds
\newcommand{\ixlrow}[3]{\texttt{#1} & #2 & #3 \\}

% one grid style for every graph; override xmin/xmax/... per plot
\pgfplotsset{sheetgrid/.style={
  axis lines=middle, grid=both, grid style={gray!45},
  xmin=-10.5, xmax=10.5, ymin=-10.5, ymax=10.5,
  xtick={-10,-8,...,10}, ytick={-10,-8,...,10},
  x=0.28cm, y=0.28cm,
  ticklabel style={font=\scriptsize, fill=white, inner sep=1pt},
  xlabel={$x$}, ylabel={$y$},
  axis line style={thick}, clip=true}}
% unit grid for small windows
\pgfplotsset{unitgrid/.style={sheetgrid,
  xmin=-5.5, xmax=5.5, ymin=-5.5, ymax=5.5,
  xtick={-5,...,5}, ytick={-5,...,5}, x=0.5cm, y=0.5cm}}
% open and closed endpoints
\tikzset{closed/.style={circle, fill, inner sep=1.6pt},
         open/.style={circle, draw, fill=white, inner sep=1.4pt}}
 then \begin{document} ... \end{document}.
\documentclass[12pt]{article}
\usepackage[margin=0.7in]{geometry}
\usepackage{amsmath,amssymb}
\usepackage{booktabs,tabularx,array}
\usepackage{enumitem}
\usepackage[T1]{fontenc}
\usepackage[utf8]{inputenc}
\usepackage{lmodern}
\usepackage{tikz}
\usepackage{pgfplots}
\pgfplotsset{compat=1.18}
\usepackage{xcolor}
\setlist{nosep,leftmargin=1.4em}
\setlength{\parindent}{0pt}
\setlength{\parskip}{6pt}
\renewcommand{\arraystretch}{1.35}
\newcolumntype{Y}{>{\raggedright\arraybackslash}X}
\pagestyle{empty}

% blanks and workspace
\newcommand{\blank}[1][2.6cm]{\rule[-2pt]{#1}{0.5pt}}
\newcommand{\lines}[1][1]{\par\vspace{2pt}\foreach \i in {1,...,#1}{\rule{0pt}{22pt}\rule{\linewidth}{0.4pt}\par}}
\newcommand{\wspart}[2]{\vspace{4pt}{\large\bfseries #1: #2}\par\vspace{-6pt}\rule{\textwidth}{0.4pt}\par}

% header: \packethead{Lesson 1-3}{Piecewise-Defined Functions}
\newcommand{\packethead}[2]{%
{\LARGE\bfseries #2}\hfill Name \blank[5cm]\\[2pt]
Algebra 2 \quad #1 (fast track) \hfill Date \blank[2.5cm] \quad Period \blank[1.2cm]\par\vspace{2pt}}

% boxed goals: \goals{learning goal}{language goal}
\newcommand{\goals}[2]{%
\fbox{\parbox{\dimexpr\textwidth-2\fboxsep-2\fboxrule}{%
\textbf{Learning goal:} #1\\[2pt]
\textbf{Language goal:} #2}}\par}

% quiz target box: \quiztarget{text}
\newcommand{\quiztarget}[1]{%
\vspace{4pt}\colorbox{gray!12}{\parbox{\dimexpr\textwidth-2\fboxsep}{\textbf{On the quiz:} #1}}\par}

% IXL table row helper: code & skill & which quiz question it feeds
\newcommand{\ixlrow}[3]{\texttt{#1} & #2 & #3 \\}

% one grid style for every graph; override xmin/xmax/... per plot
\pgfplotsset{sheetgrid/.style={
  axis lines=middle, grid=both, grid style={gray!45},
  xmin=-10.5, xmax=10.5, ymin=-10.5, ymax=10.5,
  xtick={-10,-8,...,10}, ytick={-10,-8,...,10},
  x=0.28cm, y=0.28cm,
  ticklabel style={font=\scriptsize, fill=white, inner sep=1pt},
  xlabel={$x$}, ylabel={$y$},
  axis line style={thick}, clip=true}}
% unit grid for small windows
\pgfplotsset{unitgrid/.style={sheetgrid,
  xmin=-5.5, xmax=5.5, ymin=-5.5, ymax=5.5,
  xtick={-5,...,5}, ytick={-5,...,5}, x=0.5cm, y=0.5cm}}
% open and closed endpoints
\tikzset{closed/.style={circle, fill, inner sep=1.6pt},
         open/.style={circle, draw, fill=white, inner sep=1.4pt}}
 then \begin{document} ... \end{document}.
\documentclass[12pt]{article}
\usepackage[margin=0.7in]{geometry}
\usepackage{amsmath,amssymb}
\usepackage{booktabs,tabularx,array}
\usepackage{enumitem}
\usepackage[T1]{fontenc}
\usepackage[utf8]{inputenc}
\usepackage{lmodern}
\usepackage{tikz}
\usepackage{pgfplots}
\pgfplotsset{compat=1.18}
\usepackage{xcolor}
\setlist{nosep,leftmargin=1.4em}
\setlength{\parindent}{0pt}
\setlength{\parskip}{6pt}
\renewcommand{\arraystretch}{1.35}
\newcolumntype{Y}{>{\raggedright\arraybackslash}X}
\pagestyle{empty}

% blanks and workspace
\newcommand{\blank}[1][2.6cm]{\rule[-2pt]{#1}{0.5pt}}
\newcommand{\lines}[1][1]{\par\vspace{2pt}\foreach \i in {1,...,#1}{\rule{0pt}{22pt}\rule{\linewidth}{0.4pt}\par}}
\newcommand{\wspart}[2]{\vspace{4pt}{\large\bfseries #1: #2}\par\vspace{-6pt}\rule{\textwidth}{0.4pt}\par}

% header: \packethead{Lesson 1-3}{Piecewise-Defined Functions}
\newcommand{\packethead}[2]{%
{\LARGE\bfseries #2}\hfill Name \blank[5cm]\\[2pt]
Algebra 2 \quad #1 (fast track) \hfill Date \blank[2.5cm] \quad Period \blank[1.2cm]\par\vspace{2pt}}

% boxed goals: \goals{learning goal}{language goal}
\newcommand{\goals}[2]{%
\fbox{\parbox{\dimexpr\textwidth-2\fboxsep-2\fboxrule}{%
\textbf{Learning goal:} #1\\[2pt]
\textbf{Language goal:} #2}}\par}

% quiz target box: \quiztarget{text}
\newcommand{\quiztarget}[1]{%
\vspace{4pt}\colorbox{gray!12}{\parbox{\dimexpr\textwidth-2\fboxsep}{\textbf{On the quiz:} #1}}\par}

% IXL table row helper: code & skill & which quiz question it feeds
\newcommand{\ixlrow}[3]{\texttt{#1} & #2 & #3 \\}

% one grid style for every graph; override xmin/xmax/... per plot
\pgfplotsset{sheetgrid/.style={
  axis lines=middle, grid=both, grid style={gray!45},
  xmin=-10.5, xmax=10.5, ymin=-10.5, ymax=10.5,
  xtick={-10,-8,...,10}, ytick={-10,-8,...,10},
  x=0.28cm, y=0.28cm,
  ticklabel style={font=\scriptsize, fill=white, inner sep=1pt},
  xlabel={$x$}, ylabel={$y$},
  axis line style={thick}, clip=true}}
% unit grid for small windows
\pgfplotsset{unitgrid/.style={sheetgrid,
  xmin=-5.5, xmax=5.5, ymin=-5.5, ymax=5.5,
  xtick={-5,...,5}, ytick={-5,...,5}, x=0.5cm, y=0.5cm}}
% open and closed endpoints
\tikzset{closed/.style={circle, fill, inner sep=1.6pt},
         open/.style={circle, draw, fill=white, inner sep=1.4pt}}

\usepgfplotslibrary{fillbetween}
\begin{document}

\packethead{Lesson 2-7}{Linear-Quadratic Systems}
\goals{I can solve a linear-quadratic system by substitution and shade the solution region of a linear-quadratic system of inequalities.}{I can explain why my answer solves the system using the words substitute, intersection, boundary, and region.}

\quiztarget{(1) Count the solutions of a line and a parabola (full Savvas quiz only). (2) Substitute to solve a system whose two equations both start ``$y=$'' (bonus; Example 3's first step). (3) Find how far a ball travels before it lands on a hill: set the path equal to the hill and solve (Example 3). (4) Shade the solution region of one quadratic and one linear inequality (Example 4). (5) Solve with a calculator's intersection feature (full Savvas quiz only).}

\wspart{Example 3}{Applying a Linear-Quadratic System}
Andrew kicks a ball up a hill for his dog, Laika, to chase. The hill is modeled by a line through the origin, $y=0.5x$. The path of the ball is modeled by the quadratic function $y=-0.5x^2+6x$. How far does the ball travel horizontally? How far must Laika run up the hill to catch it?

\begin{minipage}[t]{0.47\textwidth}
\vspace{0pt}\raggedright
\textbf{1.} The ball lands where the path meets the hill: an \textbf{intersection}. Both equations give $y$, so \textbf{substitute} $0.5x$ for $y$:
\[
\begin{aligned}
0.5x&=-0.5x^2+6x\\
0&=-0.5x^2+5.5x\\
0&=-0.5x(x-11)\\
x=0 \quad&\text{or}\quad x=11
\end{aligned}
\]
\textbf{2.} $x=0$ is where Andrew kicks the ball. $x=11$ is where it lands on the hill, at \mbox{$y=0.5(11)=5.5$}. \textbf{The ball travels 11 ft horizontally.}

\textbf{3.} Laika runs along the hill: the hypotenuse of a right triangle with legs 11 ft and 5.5 ft.
\[\sqrt{11^2+5.5^2}=\sqrt{151.25}\approx 12.3\ \text{ft}\]
\textbf{Laika runs about 12.3 ft.} (Answer \emph{every} part of the question.)
\end{minipage}\hfill
\begin{minipage}[t]{0.51\textwidth}
\vspace{0pt}
% TE window: first quadrant to x=27.5, y=20, ticks and grid every 2.5 ft (the TE's grid unit: "scaled by 2.5 feet"),
% x labels O,5,...,25 and y labels 5,10,15; curves arrowed. Half a unit of padding keeps the origin dot unclipped.
\begin{tikzpicture}
\begin{axis}[sheetgrid, xmin=-0.5, xmax=28, ymin=-0.5, ymax=20.5,
  xtick={0,2.5,...,27.5}, ytick={0,2.5,...,20},
  xticklabels={$O$,,5,,10,,15,,20,,25,}, yticklabels={,,5,,10,,15,,}, hide obscured x ticks=false,
  x=0.26cm, y=0.3cm,
  xlabel={Horizontal distance (ft)}, ylabel={Height (ft)},
  xlabel style={at={(axis description cs:0.5,-0.1)}, anchor=north, font=\scriptsize},
  ylabel style={at={(axis description cs:-0.08,0.5)}, rotate=90, anchor=south, font=\scriptsize}]
\addplot[very thick, red, ->, domain=0:11, samples=80] {-0.5*x^2+6*x};
\addplot[very thick, blue, ->, domain=0:27.5, samples=2] {0.5*x};
\addplot[thick, dashed] coordinates {(11,0) (11,5.5)};
\node[closed] at (axis cs:0,0) {};
\node[closed] at (axis cs:11,5.5) {};
\node[closed] at (axis cs:6,18) {};
\node[font=\scriptsize, fill=white, inner sep=1pt, anchor=south] at (axis cs:6,18.4) {$(6,18)$};
\node[font=\scriptsize, fill=white, inner sep=1pt, anchor=north west] at (axis cs:11.4,5.0) {$(11,5.5)$};
\node[font=\scriptsize, red, fill=white, inner sep=1pt, anchor=west] at (axis cs:9.8,16.5) {$y=-0.5x^2+6x$};
\node[font=\scriptsize, blue, fill=white, inner sep=1pt, anchor=south east] at (axis cs:25,13.6) {$y=0.5x$};
\node[font=\scriptsize, fill=white, inner sep=1pt, anchor=south] at (axis cs:5.5,0.3) {11 ft};
\node[font=\scriptsize, fill=white, inner sep=1pt, anchor=east] at (axis cs:10.7,2.6) {5.5 ft};
\end{axis}
\end{tikzpicture}
\end{minipage}

\textbf{Say it:} I \blank[2.6cm] $0.5x$ for $y$ because at the \blank[3cm] the path and the hill have the same $y$-value. The solution $x=11$ means \dots\lines[2]

\newpage

\wspart{Try It 3}{Revenue and cost}
Revenue for the high school band concert is given by the function $y=-30x^2+250x$, where $x$ is the ticket price, in dollars. The cost of the concert is given by the function $y=490-30x$. At what ticket price will the band make enough revenue to cover their costs?

\begin{minipage}[t]{0.48\textwidth}
\vspace{0pt}
\textbf{a.} Hint: set revenue equal to cost, then solve.\lines[4]
\end{minipage}\hfill
\begin{minipage}[t]{0.48\textwidth}
\vspace{0pt}
\textbf{b.} Hint: test a price between your two answers.\lines[4]
\end{minipage}

The band covers its costs for ticket prices from \blank[2cm] to \blank[2cm] dollars.

% R3: the lesson's only DOK 3 (Practice & Problem Solving 29, SEE EXAMPLE 3), verbatim; R5: rehearses Savvas quiz item 3 (the ball on the hill); the TE prints the path in the illustration callout.
\wspart{Practice}{Bounce or splash? (DOK 3, Practice 29)}
A golfer accidentally hits a ball toward a water hazard that is downhill from current position on the fairway. The hill can be modeled by a line through the origin with slope $-\tfrac18$. The path of the ball can be modeled by the function $y=-\tfrac{1}{100}x^2+\tfrac32x$.

\textbf{Part A} If the golfer stands at the origin, and the water hazard is 180 yd away, will the golfer's ball bounce or splash?

Hill: $y=\blank[2cm]$ \qquad Substitute and solve:\lines[4]

\textbf{Part B} How far did the ball land from the edge of the water hazard?\lines[2]

\textbf{Part C} Does it matter whether you measure the 180 yd horizontally or along the hill? Explain.\lines[2]

\newpage

\wspart{Example 4}{Solve a Linear-Quadratic System of Inequalities}
How can you solve this system of inequalities? \quad $y<-2x^2+12x-10$ \quad and \quad $4x+y>4$

\begin{minipage}[t]{0.52\textwidth}
\vspace{0pt}
\textbf{Step 1} Graph each \textbf{boundary}. ``$<$'' and ``$>$'': \textbf{dashed}, the boundary is not part of the solution.
\begin{itemize}
\item Parabola $y=-2x^2+12x-10$: vertex $(3,8)$, $x$-intercepts $1$ and $5$, $y$-intercept $-10$. Test $(0,0)$: $0<-10$? False. Shade the side \emph{without} the origin: below the parabola.
\item Line: $4x+y>4$ is $y>-4x+4$ (slope $-4$, $y$-intercept $4$). Test $(0,0)$: $0>4$? False. Shade above the line.
\end{itemize}
\textbf{Step 2} The solution \textbf{region} is where the shadings overlap (darkest). The boundaries meet where $-2x^2+12x-10=-4x+4$:
$x^2-8x+7=0$, $(x-1)(x-7)=0$, at $(1,0)$ and $(7,-24)$.

\textbf{Check a point in the region,} $(2,0)$:\\
$0<-2(2)^2+12(2)-10$, \ $0<6$ \checkmark\\
$4(2)+0>4$, \ $8>4$ \checkmark
\end{minipage}\hfill
\begin{minipage}[t]{0.45\textwidth}
\vspace{0pt}
\begin{tikzpicture}
% TE window x=[-5,10], y=[-45,25], grid 1 by 5, x labels -4,...,8, y labels 20,10,-10,...,-40; curve ends arrowed.
\begin{axis}[sheetgrid, xmin=-5.5, xmax=10.5, ymin=-45.5, ymax=25.5,
  xtick={-4,-2,2,4,6,8}, ytick={-40,-30,-20,-10,10,20},
  minor xtick={-5,...,10}, minor ytick={-45,-40,...,25}, minor grid style={gray!25},
  x=0.5cm, y=0.1cm]
\addplot[name path=par, very thick, red, dashed, <->, domain=-2.1:8.1, samples=120] {-2*x^2+12*x-10};
\addplot[name path=lin, very thick, blue, dashed, <->, domain=-5.2:10.2, samples=2] {-4*x+4};
\path[name path=bot] (axis cs:-2.1,-45.5) -- (axis cs:8.1,-45.5);
\path[name path=top] (axis cs:-5.5,25.5) -- (axis cs:10.5,25.5);
\addplot[yellow!45, opacity=0.6] fill between[of=par and bot];
\addplot[blue!30, opacity=0.45] fill between[of=lin and top];
\addplot[green!55!black, opacity=0.55] fill between[of=par and lin, soft clip={domain=1:7}];
\node[closed, red] at (axis cs:2,0) {};
\node[font=\scriptsize, fill=white, inner sep=1pt, anchor=south west] at (axis cs:2.2,0.8) {$(2,0)$};
\node[font=\scriptsize, red, fill=white, inner sep=1pt] at (axis cs:6.4,14) {$y=-2x^2+12x-10$};
\node[font=\scriptsize, blue, fill=white, inner sep=1pt, anchor=north west] at (axis cs:-5.2,9) {$y=-4x+4$};
\end{axis}
\end{tikzpicture}
\end{minipage}

\textbf{Say it:} The point $(2,0)$ is in the solution \blank[2.2cm] because \dots\ Each \blank[2.4cm] is dashed because \dots\lines[1]

\wspart{Try It 4}{Shade the solution region}
\begin{minipage}[t]{0.47\textwidth}
\vspace{0pt}
Solve the system of inequalities $y>x^2+6x-12$, $3x-y\geq-8$ using shading.
\begin{enumerate}[label=\textbf{\arabic*.}, itemsep=3pt]
\item Graph $y=x^2+6x-12$, dashed or solid?\\ Vertex: \blank[1.8cm] \ $y$-intercept: \blank[1.1cm]
\item Solve $3x-y\geq-8$ for $y$ (dividing by $-1$ flips the sign), dashed or solid?\\ $y$ \blank[3cm]
\item Test $(0,0)$ in each inequality and shade.
\item Darken the region where the shadings overlap.
\item Check a point in your region in \emph{both} original inequalities:\lines[2]
\end{enumerate}
\end{minipage}\hfill
\begin{minipage}[t]{0.50\textwidth}
\vspace{0pt}
\begin{center}
\begin{tikzpicture}
% Blank student grid (spec section 5: the window the Try It needs). It holds the vertex (-3,-21) and both boundary
% crossings (about (-6.2,-10.6) and (3.2,17.6)); the TE answer graph's window (y up to 2.5) cuts off the top of the region.
% Grid every 1 in x and every 2 in y with no gaps.
\begin{axis}[sheetgrid, xmin=-10, xmax=6, ymin=-24, ymax=20,
  xtick={-10,-8,...,6}, ytick={-24,-20,...,20},
  minor x tick num=1, minor y tick num=1, minor grid style={gray!30},
  x=0.5cm, y=0.17cm]
\end{axis}
\end{tikzpicture}
\end{center}
\end{minipage}

\newpage

\wspart{IXL}{Practice for the quiz}
Go to IXL, Algebra 2. Type the code in the search box. Do the starred skills first.

\vspace{2pt}
\begin{tabularx}{\textwidth}{l Y Y}
\toprule
\textbf{Code} & \textbf{Skill} & \textbf{Feeds} \\
\midrule
\ixlrow{HVZ $\bigstar$}{Solve a linear-quadratic system of equations: lines and parabolas}{Quiz 2 (bonus) and 3: substitute, solve the quadratic, find both points (Example 3)}
\ixlrow{VQU $\bigstar$}{Solve a linear-quadratic system of equations by graphing: lines and parabolas}{Quiz 4: graph the parabola and the line (the boundaries); Quiz 2 by graphing; Quiz 1 (full Savvas quiz only)}
\ixlrow{U5D $\bigstar$}{Solve systems of linear inequalities by graphing}{Quiz 4: dashed or solid, test a point, shade the overlap (Example 4)}
\ixlrow{BW5}{Solve a system of equations using substitution}{Warm-up: the same move with two lines (Lesson 1-6)}
\ixlrow{M69}{Solve a system of equations by graphing}{Warm-up: read where two graphs meet (Lesson 1-5)}
\bottomrule
\end{tabularx}

\vspace{4pt}
\textbf{SmartScore 70 or above on each starred skill replaces your quiz score.}

% Exit: Additional Practice 5 from the TE (2-7 Additional Practice), the move Savvas quiz item 2 (the bonus) asks for.
\wspart{Exit question}{Substitute and solve}
\begin{minipage}[t]{0.40\textwidth}
\vspace{0pt}
Use substitution to solve the system of equations.
\[
\begin{cases} y=-x^2+x+12\\ y=2x-8\end{cases}
\]
Solutions: \blank[1.8cm] and \blank[1.8cm]
\end{minipage}\hfill
\begin{minipage}[t]{0.56\textwidth}
\vspace{0pt}
\lines[4]
\end{minipage}

\textbf{Say it:} The two solutions are the points of \blank[3cm] of the line and the parabola. I checked one by \dots\lines[1]

\vspace{4pt}
\textbf{Before the quiz, finish these sentences.}

To solve a linear-quadratic system, I \blank[3cm] one expression for $y$ into the other equation and solve the quadratic.

The solutions are the points of \blank[3.2cm] of the line and the parabola: there can be 0, 1, or 2 of them.

A dashed \blank[2.6cm] is not part of the solution; the solution \blank[2.2cm] is where the shadings overlap.

\end{document}
